\documentclass[11pt]{article}
\usepackage{amsmath}
\title{Asymmetric gap margins and the deadline for one formation observation}
\begin{document}
\maketitle

Clearance and link margin of the same total width are not always
interchangeable. In the worked model, both initial gaps individually
allow the same advance, but the two require different early controls.
A one-bit report about the gap has a sharp receipt deadline even
when the common prefix is already moving.

\paragraph{Known margins.}
Use the two-robot rest-to-rest translation law: accelerations are
bounded by \(a,A\) with \(0<a\le A\), speeds by \(W,V\) with
\(0<V\le W\), and both advances begin at zero and end at the same
\(D\) at rest. With \(y=x_2-x_1\), preserve
\(-\ell\le y\le u\) throughout, with \(\ell,u\ge0\). In this
parameterization the two tolerances are relative to the actual
starting separation. Set
\[
 \tau=V/a,\quad H=T-2\tau,\quad E=VT-V^2/a,\quad
 r=(V^2/a-V^2/A)/2,\qquad
 B(t;a,v)=\int_0^t\min(a s,a(t-s),v)\,ds .
\]
For \(T\le2\tau\) the common triangular profile attains \(aT^2/4\)
and robot 1 cannot do better. Otherwise the exact maximum is
\[
       E+\min\bigl(B(H;a,W-V),\,\min(u,r)+\min(\ell,r)\bigr).
\]
At time \(\tau\), robot 2 cannot advance past
\(a\tau^2/2+\min(u,r)\), by the gap bound and its own speed and
acceleration envelope. Reverse time to get the lower bound
\(D-a\tau^2/2-\min(\ell,r)\) for its advance at \(T-\tau\).
In between it contributes at most \(VH\). The remaining upper bound
is robot 1's standalone capacity.
For sharpness choose a start lead \(p\le\min(u,r)\) and end lag
\(n\le\min(\ell,r)\) whose sum is the required extra advance.
These are nonnegative. Keep robot 1's ramps of length \(\tau\);
robot 2 reaches \(V\) at time \(\tau-2p/V\), and uses a reversed
ramp with \(n\) at the end. In the middle it travels at \(V\)
while robot 1 adds a scaled capped triangular profile of area
\(p+n\). The relative position increases to \(p\), decreases
to \(-n\), then returns to zero. Every speed and acceleration
bound holds. For \(a=V=1,A=4,W=3,T=4\), bands
\([-1/4,1/4]\) and \([0,1/2]\) have the same width yet permit
advances \(3.5\) and \(3.375\) respectively.

\paragraph{Unknown initial gap, known delivery time.}
Keep those actuator bounds, deadline 4, and physical gap band
\([1,8/5]\). The initial gap is either \(6/5\) or \(7/5\),
giving relative bands \(A=[-1/5,2/5]\) and \(B=[-2/5,1/5]\).
The command source cannot otherwise tell the cases apart before
a truthful bit arrives simultaneously for both robots. Actual
motion is deterministic given commands, with relative advances
but no other gap observation; no clocks, dynamics or state
measurements reveal the initial separation before the bit.
Both cases separately permit target \(D=143/40\). With no bit,
both advances must coincide in the two worlds, so the intersection
band limits the maximum to \(17/5\).
For that target the latest receipt time is exactly
\[
                              c=1-\sqrt{7}/4.
\]

To prove the cutoff, each case must saturate the two end-margin
bounds and middle speed bound in the known-margin calculation.
In case A, \(x_2(1)=7/8\) forces
\(v_2(t)=\min(4t,1)\) throughout \([0,1]\).
In case B, \(x_2(1)=7/10\) forces \(x_1(1)=1/2\), hence
\(v_1(t)=t\) throughout \([0,1]\).
Equality in each position integral forces equality of the
velocity envelope, since the velocities are continuous.
For a shared prefix through \(s\in[1/4,1]\), case B consequently
has \(y(s)=s-1/8-s^2/2\), with \(y'(s)=1-s\).
Relative deceleration is at most \(5\); before relative
velocity can be nonpositive the gap must reach at least
\[
             s-1/8-s^2/2+(1-s)^2/10
             =3/8-(2/5)(1-s)^2 .
\]
The time \(s+(1-s)/5\) is still no later than 1.
To stay in B this peak must not exceed \(1/5\), which gives
\(s\le c\). A receipt at or after 1 already violates the
bound on the shared prefix at 1; earlier than \(1/4\) is not
excluded by the necessary condition.

\paragraph{Attaining the cutoff.}
Share \(v_1=t\) and \(v_2=\min(4t,1)\) through \(c\).
Case A then preserves both initial profiles until time 1.
Case B brakes robot 2 at 4 until \(b=1-\sqrt7/5\), when it has
the same speed as robot 1, then sets \(v_2=t\) through time 1.
The gap in B hits \(1/5\) at \(b\) and does not exceed it.
For the middle interval use \(v_2=1\) and
\(v_1=1+(23/40)\min(t-1,3-t)\) in both cases.
The final reversed ramps use lag \(1/5\) in A and \(3/8\) in B.
They reach \(D\) and stop together. Earlier report receipt can
use this same policy and wait to branch at \(c\).
The replay verifies every piecewise gap extremum, including
inside intervals, using exact arithmetic in \(\mathbb Q(\sqrt7)\).

A direct sampled shared-prefix optimizer matches the target at
receipt 0.33 and falls short at 0.35; the exact threshold follows
from the proof rather than those samples. Compositional barrier
certificates address continuous safety and connectivity; the
additional question here is when a private observation must arrive
to meet a joint deadline. The bit is an ideal input. This does
not compare implementations or establish existence of a real
gap sensor at the cut. The result is scoped R6, not a field or
independent mission result.
\end{document}
